% English draft based on sections/01_preliminaries_ja.tex
\section{Preliminaries: VED Vocabulary in the Intelligence Layer}

This chapter fixes the vocabulary needed for the rest of the paper. The aim is not to import all of VED or IFGT, but to isolate the minimum structure required to describe inference, stopping, projection, and emergent stopping layers.

The IFGT vocabulary used here follows the projected-uses stance adopted in the consciousness paper. Symbols such as $I$, $J$, $\Phi$, and projection are not direct identifications with the full IFGT field. They are effective coordinates projected into the intelligence-layer observation window.

\subsection{Methodological Fixing: Stopping, Projection, and Consensus}

Stopping is not truth. In this paper, a stopping point is a point at which inference can be operated as provisionally stable. It may later be revised. Accordingly, the closure degree $C$ should be read as a stability measure, not as a truth measure.

The threshold mobility $\mu_\theta$ is not mere flexibility. It is the ability to reorganize where a stopping operator $E$ can be installed. Sapiens-like intelligence is not characterized by the absence of stopping. It is characterized by having stopping while also being able to relocate stopping.

Projection and consensus are read as operations of stopping structure. Attachments to ideals, gods, communities, money, scientific theories, and other loci can be structurally interpreted as installations of stopping operators. This paper does not enter psychological or clinical accounts of projection, but it keeps the structural relation between projection and stopping explicit.

The social blockchain discussed later is not blockchain as version-control technology. Its center is not Version Control but Consensus: the sharing and synchronization of stopping structures across agents.

Science and religion are different implementations of stopping-operation systems. Religion is strong in long-term maintenance of stopping structures; science institutionalizes the revision of stopping structures. The increase of inferential capacity does not itself guarantee stability. Rather, more advanced inference requires greater stopping cost. In this sense, religion and science are distinct operating forms that make inference sustainable under finite bodies and finite resources.

\subsection{VED Basic Form}

VED begins from difference. Where there is difference, there is a gradient; where there is a gradient, there is flow; where flow is constrained, quasi-closure and attractors can emerge. In the intelligence layer, this grammar appears as:
\[
\text{difference} \;\longrightarrow\; F(L_F) \;\longrightarrow\; \text{non-closure} \;\longrightarrow\; \text{externalized stopping}.
\]

The closure degree $C$ measures how close a state is to operable stability. The information-field drive can be written schematically as a function of non-closure:
\[
I = f(1-C).
\]
As $C$ approaches an unattainable $C_{\max}$, the system does not arrive at absolute closure. Instead, the drive toward closure becomes singular or unstable.

\subsection{Movement Region and Threshold Mobility}

Let $\Omega_\theta$ denote the movement region available under a constraint parameter $\theta$. We distinguish two movement types:
\begin{itemize}
\item \textbf{(A) movement within a fixed region}: the state moves inside a given $\Omega_\theta$;
\item \textbf{(B) reconfiguration of the region}: $\theta$ itself changes, thereby changing $\Omega_\theta$.
\end{itemize}

The non-closure theorem in \S{}3 applies to (A), movement within fixed $\theta$. The emergence mechanism in \S{}5 requires (B), region reconfiguration. This distinction is essential. It allows the paper to say that inference cannot stop itself within a fixed region while also explaining how stopping can emerge when the region of possible movement is reorganized.

\subsection{Emergence}

This paper uses emergence in the VED sense:
\[
U = \Phi(L,H_{ij}),
\]
where lower-layer motion $L$ and external mutual constraint $H_{ij}$ give rise to an upper-layer attractor $U$. Emergence is not a vague increase of complexity. It is the formation of an attractor from non-closed motion under mutual constraint.

For Intelligence Part II, the relevant substitution is:
\[
L_E = \Phi(L_R,H_{ij}).
\]
The emergent stopping layer $L_E$ is thus not deduced from the inference layer alone. It arises from the recursive motion of the log layer $L_R$ under external mutual constraint.

\subsection{Variables}

\begin{longtable}{@{}p{0.22\linewidth}p{0.68\linewidth}@{}}
\toprule
Symbol & Role \\
\midrule
\endfirsthead
\toprule
Symbol & Role \\
\midrule
\endhead
$L_F$ & inference layer; the layer on which the inference operator $F$ acts \\
$L_R$ & log layer; recursive reference to traces left by inference \\
$L_E$ & emergent stopping layer; stopping attractor produced from $L_R$ and $H_{ij}$ \\
$H_{ij}$ & mutual constraint between agents, or between agent and environment \\
$\mu_\theta$ & threshold mobility; ability to reorganize the possible installation region of $E$ \\
$E$ & external stopping operator acting through $L_E$ \\
$\Phi_{\text{ext}}$ & external reference field as it appears from the perspective of an individual \\
$\Phi_{\text{net}}$ & inter-individual network of $L_E$ layers \\
\bottomrule
\end{longtable}

\bigskip
